1
A Method to Align Goals and Business Processes
Renata Guizzardi, Ariane Nunes Reis
Ontology & Conceptual Modeling Research Group
Federal University of Espírito Santo
Vitória, Brazil
rguizzardi@inf.ufes.br, arianenreis@gmail.com
Abstract. Business Process Modeling (BPM) has been for a number of years in
the spotlight of research and practice, aiming at providing organizations with
conceptual modeling-based representations of the flow of activities that generate
its main products and services. It is essential that such flow of activities is engineered in a way to satisfy the organization’s goals. However, the work on BPM
still makes shy use of goal modeling and the relation between goals and processes
is often neglected. In this paper, we propose a method that supports the analyst
in identifying which activities in a business process satisfy the organization’s
goals. Moreover, our method allows reasoning regarding the impact of each of
these activities in the satisfaction of the strategic (i.e. top) goals of the organization. The results of this analysis may lead to reengineering, and grant the analyst
with the means to design higher quality BPMs. Besides describing the method,
this paper presents a preliminary evaluation of the method by the means of an
empirical study made in a controlled environment.
1 Introduction
Competitive businesses and an ever-changing market have demanded that current organizations constantly evolve. To achieve that, it becomes necessary to develop a deep
understanding of the organizational processes and systems. This motivates an increasing interest in Business Processes Modeling (BPM), the discipline concerned with explicitly capturing, by applying conceptual modeling languages, the flow of activities
that generate the main products and services offered by the organization [1].
However, modeling the flow of activities may not be enough to provide the organization with competitive advantage. It is also important to grasp if the current activities
and business processes are in line with the goals of the organization. This idea is supported by Rosemann and vom Brocke [2], who claim that strategic alignment is one of
the six core elements of BP Management. Although goal modeling is often supported
by BPM platforms, it has been regarded by BP practitioners as of secondary importance, and little is explored regarding the relation among processes and goals.
Understanding if and how processes achieve the operational and strategic goals of
the organization may guide the decision regarding which activities and processes
should become priorities. This may be realized by trying to align (map) activities and
business processes to the organization's goals. Moreover, aligning processes and goals 
may help the analyst to understand if there is any goal being neglected. If a goal is not
aligned to any activity or business process, then one of the following may be happening:
i) the business process is poorly modeled and an activity should be added to make explicit something that is already done in practice; ii) the business process needs to be
reengineered to accommodate one or more new activities, so that the goal is accomplished; or iii) the goal is no longer important to the organization and should be removed
from the goal model.
Many organizations today develop information systems based on their modeled business processes [3]. Thus, aligning goals and business processes may also assist in the
development of systems that are more in line with the organization's goals.
Although, recently, some works have focused on the alignment of goals and processes[4,5,6,7], a systematic method to support goal and process alignment is still missing. This paper addresses this gap, by describing the results of an ongoing effort towards
the alignment of goals and business processes. In a previous work, we have theoretically addressed this topic [8], while this paper focuses on a step-by-step method to support goals and business processes alignment. This method aims at modeling which particular activities or processes achieve the goals of the organization, also providing a
reasoning mechanism to verify the impact of the execution (or non-execution) of these
activities and processes have on goal satisfaction. Moreover, the method allows one to
find inconsistencies in the BP model, supporting the analyst in building models of
higher quality and in designing business models more in line with organizational strategies. In this context, goals play the role of drivers of business process improvement.
The remaining of this paper is organized as follows: Section 2 presents the applied
goal modeling and BPM languages; Section 3 describes the proposed method, illustrating its steps with the use of a running example; Section 4 discusses a preliminary evaluation conducted by the means of an empirical study; Section 5 compares the method
with some related works; and finally, Section 6 concludes this paper.
2 The Adopted Modeling Languages
2.1 Goal Modeling Language and Reasoning Mechanism
Languages and solutions supporting BPM still make shy use of goal and strategic modeling. This is the case, for example, of the ARIS platform, one of the most used BPM
solution in industry. While supporting goal modeling, ARIS proposes a modeling language with very low expressivity, which basically allows relating macro-processes to
the leaves of a goal tree, without however distinguishing different kinds of relations
between goals and processes and with no reasoning support. Works developed in the
context of GORE (Goal-Oriented Requirements Engineering), on the other hand, provide much more powerful modeling languages such as the ones adopted by i* [9] and
KAOS [10], for example. Aiming at profiting from this earlier work, we here adopt the
Tropos language [11], an i* dialect.
In Tropos, goals are modeled in the perspective of an actor, and often lead to treestructures, where top goals are refined into lower-level ones, mainly by the means of 
AND and OR decompositions. As one can grasp by intuition, when a top goal is decomposed in two sub-goals by the means of an AND-decomposition, only the satisfaction of the two sub-goals lead to the satisfaction of the top one. If however, a top goal
is OR-decomposed into two sub-goals, it is sufficient that only one of them is satisfied
in order to satisfy the top goal. Figures 2, 4 and 5 are examples of Tropos models.
Giorgini et al. [12] propose a Tropos-based quantitative approach for verifying goal
satisfaction, based on the propagation of values in the goal tree. This approach allows
two kinds of propagation, forward and backward propagation. In the forward propagation, initial satisfaction values are attributed to the leaf goals and are propagated to the
root goals, thus determining how much evidence there is to the satisfaction of the root
goal, given the probability of satisfaction of the leaf goals; while in backward propagation, a desired final satisfaction value is given to the root goal, and then propagated
from the root to the leaf goals, so as to determine how much evidence regarding the leaf
goals satisfaction is need in order to obtain the desired evidence of satisfaction for the
root goal.
To propagate the satisfaction values, it is necessary to take into consideration the
labels of the relationships between the goals, as follows:
 AND-decomposition: the decomposed goal is satisfied only with the complete satisfaction of the subgoals. However, even if the evidence of all the subgoals’ satisfaction is partial, this evidence propagates to the decomposed goal.
 OR-decomposition: the decomposed goal is satisfied if one or more subgoals are
completely satisfied. However, even if the evidence of one of the subgoals satisfaction is partial, this evidence propagates to the decomposed goal.
 Contribution: the contributing goal’s satisfaction propagates to the contributed
goal, based on a particular contribution (impact) weight. The reasoning approach
supports positive and negative contribution weights, both for goal satisfaction and
denial. For reasons of space limitation, in this paper we only consider positive contributions, for the case of goal satisfaction.
2.2 Business Process Modeling Notation
For modeling BPs, we adopt BPMN [13], which has become a well-accepted standard, being implemented in many BPM tools and largely adopted, both by researchers
and practitioners. BPMN allows the flow of activities to be modeled, by depicting the
order in which activities occur. The BPMN elements used in this paper are:
 pool: represents a participant (actor) in the process, delimiting a space where
the activities performed by that participant are modeled;
 task (rectangle): represents a process activity;
 start event (circle): triggers the beginning of the process;
 end event (circle): determines the process is ended;
 intermediate event (thick-boarded circle): models an event that occurs in the
middle of the process, usually leading to some decision and altering the flow
of activities.
 exclusive gateway (diamond): placed in the beginning of a fork of activities,
determines that only one of the forked paths is followed.
Figures 1 and 3 present examples of BPMN models.
3 A Systematic Method to Align Business Processes and Goals
The proposed method assumes that the organization’s goals and processes have been
previously modeled, taking both models as entries. While valuing approaches that start
with goal modeling and then model processes based on these goals, e.g. [5,6], we also
recognize that, in practice, many organizations have undertaken these modeling tasks
separately and today have several non-aligned goal models and BPMs. In this context,
it is possible that goal and process model have different granularity level. If that is the
case, it may be necessary to refine the models, as our method assumes these models
have the same granularity. The general process underlying the proposed method is composed of the five steps depicted in Figure 1. Aiming at illustrating each of these steps,
subsection 3.1 describes a running example, further referenced in the remaining subsections to demonstrate how each step of the alignment method is carried out.
Fig. 1. The process underlying the proposed method
3.1 Running Example
For the running example, we choose a scenario of conference paper reviewing, for
being a well-known and simple enough example to enable the method illustration. Figures 2 and 3 present this scenario’s goal and process model respectively.
Fig.2. Conference paper reviewing goal model
Fig. 3. Conference paper reviewing BPMN model
From now on, we highlight the names of goals and activities from these models,
using for them a special font, to facilitate reading.
3.2 Classify process paths
In the first step, the analyst classifies the process paths as one of the following three
types:
 Main path: expected path followed by the process;
 Secondary path: forking path whose end leads back to the main path;
 Alternative path: forking path that does not return to the main path, hence
providing an alternative ending to the process.
This classification is important because, for a given process, some goals may be only
satisfied by activities in the main path that, in a certain process execution (i.e. instance), 
are never actually reached. The opposite situation is also problematic, when goals are
only satisfied by activities belonging to secondary or alternative paths.
In the conference paper review process of Fig.3, the path containing the activity
named notify track cancelation to conference chairs is an alternative path, and the
path containing the review papers with missing reviews activity is a secondary path.
All activities not belonging to these two paths compose the main process path.
3.3 Assign activities and sub-processes to goals
In this step, the analyst takes each leaf goal from the given goal model and assign it
to an activity or a sub-process that leads to this goal accomplishment. We call this goal
and activity alignment. If the goal is aligned to an activity, it is classified as an activity
goal; if on the other hand, it is aligned to the process as a whole or to one of its subprocesses, it is classified as a process goal. This classification is done according to the
taxonomy provided in [8] and it provides for the analyst the information if such a goal
is achieved after the execution of solely one activity or if several activities should be
taken into account to determine if a goal is successfully accomplished.
Given the aforementioned alignment, the analyst is able to add plans in the goal
model, each one representing an activity of the existing process. Plans are linked to
goals via the Tropos positive contribution relation, stating that a given plan influences
positively the satisfaction of the goal it is linked to. Figure 4 presents the resulting goal
model, after the alignment and classification. Process goals are depicted in light grey
while activity goals are depicted in white.
Fig. 4. Goal model after activity alignment
As can be noted by the alignment of the distribute all papers to PC members
activity, it is possible for an activity to contribute to the satisfaction of more than one
goal. Many alignments are straightforward, in a way that each activity coming from the
BPMN model is designed as a Tropos plan with the same name, in the Tropos model.
Some alignments, however, may require some design adaptation, as in the case two or
more activities from the process compose a sub-processes aligned to a goal, which normally leads to the creation of a super-plan representing the sub-process. This is the case
of the have papers reviewed by experts in the field goal, for which we created the
involve PC members plan, which is then AND-decomposed into select PC Members and invite PC Members (the two activities coming from the process). This may
also indicate to the analyst that the process model should be changed, to substitute given
activities by a sub-process1
, so as to create a simpler and easier-to-read process. Another case for design adaptation is related to the path classification discussed in section
3.2. When a goal is accomplished both in the main path and in an alternative or secondary path, the plans representing the activities that achieve such goal should be related
via OR-decomposition. This makes sense because only one of the activities will be
executed at each process instance, depending if the secondary or alternative path was
followed, or if the process executed its main path. This is the case of the have reviews
on time goal. This goal is accomplished both if all papers are reviewed by PC Members
(the top path in the identified secondary path) or if the PC Chair reviews the missing
papers himself (the bottom path in this same path). To model that, an auxiliary plan
named review all papers (in time) was created and then decomposed into all papers
are reviewed by PC Members and review missing papers, only this last one coming straight from the process model.
In Tropos, there are both positive and negative contributions, respectively meaning
positive and negative influences on goal satisfaction. Since our method concerns process and goal alignment, for the sake of exemplifying it, we assume that all activities
in the process are executed in order to accomplish the organization’s goals, thus we
only consider positive contributions. Nevertheless, we are aware that an organization
may have conflicting goals. In that case, a process activity may have a positive impact
on the satisfaction of one goal, while undermining the accomplishment of another. In
such case, positive and negative contributions must be used and are supported by the
reasoning mechanism we apply in our method.
3.4 Verify non-aligned goals and activities
In this step, the analyst should examine the model to check: i) if each activity in the
process is aligned to at least one leaf goal in the goal model; ii) if all goals from the
goal model are satisfied by activities or sub-processes in the process; iii) if there is some
goal being neglected in case the process follows a given path. When examining the goal
and process model for our running example, we find problems when performing all
these three verifications. Let us then examine each of these cases.

1 BPM languages (as BPMN) and frameworks usually offer the option of including as activities
in the process model, subprocesses that are also modeled separately, as a modularity strategy.
A process activity is not aligned to any goal in the goal model. By looking back
to the process model of the running example, we realize that there are two activities not
aligned to any goal in the goal model: send out call for papers and notify end of
review to conference chairs. At this point, we should verify if these activities are
simply not relevant and exceeding in the present process, or if new goals satisfied by
such activities should be added to the goal model.
The send call for papers activity is essential, otherwise, nobody will find out about
the track and, thus, no paper will be submitted. For this reason, we create a new goal in
the goal tree, named attract submissions, which can be directly linked to a send call
for papers plan, representing this activity.
When performing the step of aligning goals and activities (refer to section 3.3), we
realized that the notify end of review to conference chairs activity is actually redundant, because, when we submit the list of accepted papers to the chairs, we are already
communicating about the end of the review process. We conclude that such activity is
in fact exceeding and exclude it from the process.
A goal is not satisfied by any activity or sub-process. Fig. 4 shows that the avoid
bias goal has no aligned activity. This means that such goal is never achieved during
any execution of the given process. This may be because the goal is not actually important and should, thus, be excluded from the goal model. On the other hand, it may
be a clue for the analyst that the process should be reengineered so as to include activities to accomplish such goal. In this case, we believe the latter, so the ask PC Members to perform bidding activity is added in the process, right before the distribute
all papers to PC Members activity, thus allowing PC Members to declare conflicts
and express their preference regarding papers to review.
Whenever there are missing activities, this may be the result of poor modeling, thus
adding activities contributes to enhancing the quality of the process model. In the worst
case, missing activities mean that the process must actually be reengineered. In both
cases, our method shows that goals may be the drivers to account for such missing
activities, supporting the work of the BP analyst.
A goal is neglected in a particular process path. By analyzing the model after path
classification, we realize that the keep good communication with paper authors
goal is neglected when the process takes the alternative path. In other words, in the first
forking, if there are not enough submissions, the conference chairs and the PC members
are notified, however, the authors of the submitted papers are never informed regarding
the track cancelation. In this case, again, we must add an activity in the process, right
after the notify track cancelations to conference chairs activity, named notify track
cancelation to authors.
It is important to realize that in the alternative path, many of the goals that are not
pertinent to this path are actually neglected. For example, the papers are never going to
be reviewed so the have papers reviewed by experts in the field goal and the avoid
bias goal will never be achieved. However, this results from an abnormal termination
of the process and should not lead to any process change.
Another case in which goals may be correctly achieved only in a secondary path
regards the case in which there are two alternative sub-goals (i.e. sub-goals related via
OR-decomposition). Suppose, for instance, that the conference track may have either
invited submissions or papers submitted in response to a call. In that case, besides the 
attract submissions goal, there would be an alternative goal named invite submissions. In the process, the send out Call for Papers activity would be forked with a
send invitation email to known authors activity. Each of these secondary paths
would accomplish only one of the previously mentioned goals, and this would be perfectly natural.
3.5 Attribute impact values to plans
In this step, the analyst (in collaboration with the stakeholder) should attribute
weights (i.e. impact values) for the contribution links in the goal model. Given that a
plan P contributes to a goal G, a contribution weight means how much impact P has on
G’s satisfaction. Following the probabilistic model proposed in [12], consider a goal G,
which has two contributing plans P1 and P2, whose contributions are attributed values
of +0,4 and +0,7, respectively. This means that if plan P1 is completely executed, there
is a 40% chance that goal G is achieved, while if plan P2 is completely executed, this
represents a greater impact on goal G (70% chance of achievement). Figure 5 shows
the goal model of the running example after the contribution weights attribution.
Fig.5. The goal model after impact values are attributed to goals and plans
3.6 Values Propagation
In this step, the analyst should attribute different satisfaction values to the leaf plans
of the model and simulate how these values propagate in the goal tree, impacting the
middle and root goals. Here, we adapt the proposed approach [12], by considering that
for a plan, the satisfaction value consists in the probability of plan execution. For reasons of lack of space, we exemplify the reasoning method only for contribution value
propagation in case of satisfaction. Thus, in case the plan is not executed, nothing can
be said about the denial of the goal to which the plan contributes. The attributed values
should be propagated from the leaves to the top, following the propagation rules proposed in [12] and shown in Table 1. Consider Sat(G1) as the satisfaction value of G1
(i.e. the probability that G1 is satisfied) and w, the contribution weight, explained in
section 3.5.
Table 1. Propagation rules
Contribution 𝐺2
𝑤+𝑆
→ 𝐺1
𝑆𝑎𝑡(𝐺1
) = 𝑆𝑎𝑡(𝐺2
). 𝑤
AND-decomposition (𝐺2, 𝐺3
)
𝑎𝑛𝑑
→ 𝐺1
𝑆𝑎𝑡(𝐺1
) = 𝑆𝑎𝑡(𝐺2
). 𝑆𝑎𝑡(𝐺3
)
OR-decomposition (𝐺2, 𝐺3
)
𝑜𝑟
→ 𝐺1
𝑆𝑎𝑡(𝐺1
) = 𝑆𝑎𝑡(𝐺2
) + 𝑆𝑎𝑡(𝐺3
)
− 𝑆𝑎𝑡(𝐺2
). 𝑆𝑎𝑡(𝐺3
)
Suppose that in our example, all plans referring to activities in the main path of the
process have been fully executed (thus, having satisfaction value=1), except the all papers reviewed by PC members by review deadline plan, whose satisfaction value
is 0,9. This leads the process to its secondary path, and we will suppose that the review
papers with missing review plan is also fully executed (thus, having satisfaction
value=1). Based on these initial values and using the rules of Table 1, Table 2 presents
the propagated values to the remaining plans and goals of the model.
Table 2. Propagated values to the plans and goals in the goal model
Plans Sat Plans Sat
involve PC members 1 keep PC members informed 1
review all papers 1 keep authors informed 1
keep conference chairs informed 1
Goals Sat Goals Sat
organize successful conference
track
0,6 have papers reviewed by experts
in the field
0,9
attract submissions 1 participate as little as possible in
the review process
0,9
guarantee papers have high quality
reviews
0,7 have reviews on time 1
guarantee efficiency in the review
process
0,9 with the conference chairs 1
Keep a good communication with
all actors involved
1 with the PC members 1
avoid bias 0,8 with the paper authors 1
To exemplify how the values are propagated, let us analyze the evidence of satisfaction of the guarantee efficiency in the review process goal. For that, we start from
the leaf plans that are indirectly related to this goal. The all papers reviewed by PC
members by review deadline plan (satisfaction value=0,9) and the review papers
with missing review plan (satisfaction value=1) are related via OR-decomposition to
the review all papers plan (satisfaction value=0,9+1-0,9*1=1). The review all papers
plan contributes with a weight of +1 to the have reviews on time goal (satisfaction
value=1*1=1). By analogous calculations, we arrive at the value of 0,9 to the satisfaction value of the participate as little as possible in the review process goal. And 
finally, the participate as little as possible in the review process goal and the have
reviews on time goal are related via AND-decomposition to the guarantee efficiency
in the review process goal (satisfaction value=0,9*1=0,9). In case there are two plans
contributing to a goal (e.g. in the case of the have papers reviewed by experts in
the field goal), only the maximum satisfaction value is propagated.
By testing different values assigned to the leaf plans, the process manager is able to
reason about the impact of activity execution in the organization’s goals. This may help
him set priorities and allocate resources to the different process activities, based on their
impact on goal achievement. In case the process is automated via a workflow system,
after the execution of each activity, the system can automatically update a goal model
like the one we are using, providing it with new values of goal satisfaction. This will
give the process manager the ability to monitor, in real time, how the middle and the
top goal are being achieved.
4 Preliminary Evaluation of the Method
With the aim of evaluating the proposed method, an empirical study was conducted,
with the voluntary participation of 14 students from a master course in Computer Science, at the Federal University of Espirito Santo (UFES), in Brazil. For this study, we
used real process models from a public regulatory agency. We obtained the real alignment from the employees of such agency and considered it for comparison with the
results performed by the empirical study participants. This allowed us to count how
many correct or incorrect alignments there were for each participant, as well as the
absence of correct alignments.
When quantitatively analyzing the results, we considered the following metrics:
 Efficiency metrics: a) the rate between the number of correct alignments and the
time spent performing alignments; b) the rate between the number of incorrect
alignments and the time spent performing alignments; c) the rate between the number of absent correct alignments and the time spent performing alignments;
 Effectiveness metrics: d) the rate of correct alignment and the total number of alignments; e) the rate of incorrect alignment and the total number of alignments; f) the
rate of absent correct alignment and the total number of alignments.
The study was conducted in two rounds and the participants were divided in two
groups (A and B), considering their levels of expertise regarding goal and process modeling. Group A aggregated the more experienced participants in both areas, while group
B aggregated those less experienced. Information regarding participant´s expertise was
gathered by the analysis of a profiling form filled by the participants. In the first round,
both groups performed the ad-hoc alignment, i.e. each participant had to individually
align activities and goals without the use of any particular method. In the second round,
the participants of both groups received instructions regarding the proposed method,
each one performing a new alignment based on this method. Two cases of similar complexity were chosen for the empirical study. In the first round, group A had case 1,
while group B had case 2 (in the second round, this was inverted). Moreover, to enable 
a qualitative analysis, we designed a form, in which the participants could share their
impressions and explicitly give suggestions for the improvement of the method.
For the quantitative analysis, we used boxplot, which allows the visualization of the
concentration of the collected data, excluding outliers and supporting the results’ comparison. In terms of efficiency, the result showed that the only metric in which the ad
hoc alignment performed better is the rate of the number of correct alignment and the
time (metric a). This is justifiable, given that the proposed method requires several steps
while in the ad hoc method, the alignment is directly made by intuition. Moreover, the
fact that all participants were novices in the use of the proposed method may also have
influenced the longer time spent in performing alignments. However, in terms of avoiding mistakes (metrics b and c), the proposed method performed better. In terms of effectiveness, the proposed method led to a higher number of correct alignments (metric
d) and a lower number of incorrect alignments (metric e). Nevertheless, the ad hoc
method performed better in avoiding absence of correct alignments (metric f).
Concerning the qualitative analysis, the participants practically unanimously evaluated the proposed method positively. The biggest contributions of the method, according to the participants is the existence of objective steps that allow them to: i) systematically identify the goal and activity alignments. With respect to this point, the participants emphasized the usefulness of the process path classification (see section 3.2) and
goal classification (see section 3.3) to facilitate the alignment; and ii) detect discrepancies between the goal and process models (see section 3.4).
Although the results were positive, the participants also highlighted some important
limitations: 1) they found the method very complex, being composed of too many steps
and sub-steps; and 2) they also claimed that the time dedicated for training was not
enough. This possibly explains the bad quantitative results with respect to method efficiency. Moreover it indicates that, when applied in practice, the method requires heavy
training and its success relies on the analyst’s level of experience.
5 Related Works
The most related research to our own is the GPI method [4]. GPI combines i* and
BPMN, aiming at profiting from each of these languages’ dispositions in describing the
organization’sstrategic and operational dimensions, helping to maintain the traceability
between goals and processes. Like in our work, the integration between these two languages is obtained through the i* task element (plan in Tropos). However, differently
than in our misalignment detection strategy, goal satisfaction relies on monitoring KPIs
related to the resources produced by the BPs that are aligned to the goal. If the given
process does not produce an expected KPI, then the goal is not satisfied and there is a
misalignment in the organization.
Nagel et al. [5] report on an approach to ensure consistency among goals and BPs,
based on the use of KAOS and applying model checking. In their work, the focus is on
determining the order in which activities should occur, given by logical and temporal
dependencies between goals.
Jander et al. propose Goal-oriented Process Modeling Notation (GPMN) [6], a specific language to model goal-oriented BPs. As in our approach, these authors recognize
the importance of explicitly modeling the motivations behind the execution of the processes. However, they propose a different notation instead of using existing goal-modeling languages. We prefer to reuse existing works on goal modeling, so as to build
over a big body of work previously done in this research field.
Soffer and Rolland [7] propose a method to combine intention-oriented (i.e. goal
modeling) and state-based process modeling. Regarding the modeling languages, instead of Tropos and BPMN, this work adopts two state-based modeling languages. In
general, it consists in a more formal work when compared to ours. However, we do
share some objectives, such as detecting model incompleteness.
Differently than us, both Nagel et al. [5] and Jander et al. [6] consider that the BP
model starts with goal-modeling and systematically proceeds to activities’ modeling.
Thus they do not pay too much attention to the alignment inconsistencies as described
in section 3.4. Moreover, a limitation of the works presented in [5,6,7] in comparison
to ours is not providing any reasoning mechanism regarding goal satisfaction.
6 Conclusion
This paper presents a method for aligning goals and process activities, based on the
application of state of the art techniques in goal modeling and BPM, namely Tropos
and BPMN.
Goal and process alignment may help the analyst to detect inconsistencies in the goal
and process models, thus leading him to enhance the quality of the models and, ultimately, to reengineer. The method also helps the project manager to estimate the impact
of the execution of each process activity in goal satisfaction, thus assisting him in setting priorities and making decisions regarding the time, effort and resources spent with
each activity.
We are aware that it may be hard to identify the real values of impact each activity
has on goal achievement, since these values are subjective and hard to determine. However, we believe the stakeholders will have at least an idea if the impact is low, medium
or high, allowing him to set different values and test them. For the future, we will explore qualitative reasoning mechanisms to evaluate if they are more suitable and practical.
Besides describing the method, the paper presents a preliminary evaluation made
with an empirical study. We acknowledge the limitations of this evaluation, especially
regarding the low number of subjects. However, we believe that it provides a good
indication and a basis for a full experimental evaluation, aimed for the future. Moreover, we aim at performing real case studies, expecting that results from real environments will help us understand the viability of the method, also helping us improve it.
Acknowledgement. This work is partially supported by CAPES/CNPq (grant number
402991/2012-5) and CNPq (grant numbers 461777/2014-2 and 485368/2013-7). The authors are 
also grateful to Giancarlo Guizzardi, João Paulo Andrade Almeida and Mateus C. Barcellos Costa
for the fruitful discussion regarding the proposed method.
References
1. Davies, I., Green, P., Rosemann, M., Indulska, M., Gallo, S.: How do Practitioners Use
Conceptual Modeling in Practice? Data & Knowledge Engineering, 58 (2006), 358--380.
2. Rosemann, M., vom Brocke, J.: The six core elements of business process management. In
Rosemann, M., vom Brocke, J. (eds.) Handbook on Business Process Management, 1,
Springer-Verlag, Berlin (2015), 105--122.
3. Dumas, M., van der Aalst, W.M.P., ter Hofstede, A.H.M. (eds.): Process Aware Information
Systems: Bridging People and Software Through Process Technology. John Wiley & Sons,
New Jersey (2005)
4. Sousa, H., Leite, J.: Modeling Organizational Alignment. In: 33rd International Conference
on Conceptual Modeling (2014), 407--414.
5. Nagel, B., Gerth, C., Engels, G., Post, J.: Ensuring Consistency Among Business Goals and
Business Process Models. In: 17th IEEE International Enterprise Distributed Object Computing Conference. (2013), 17--26.
6. Jander, K., Braubach, L., Pokahr, A., Lamersdorf, W., Wack, K.: Goal-oriented Processes
with GPMN. In: International Journal on Artificial Intelligence Tools, 20, 6. (2011), 1021—
1041
7. Soffer, P., Rolland, C.: Combining Intention-Oriented and State-Based Process Modeling,
In: 24th International Conference on Conceptual Modeling (2005), 47--62.
8. Cardoso, E., Guizzardi, R., Almeida, J.P.: Aligning Objectives and Business Process Models: A Case Study in the Health Care Industry. International Journal of Business Process
Integration and Management, 5, 2 (2011) , 44--158
9. Yu, E., Giorgini, P., Maiden, N., Mylopoulos, J. (eds.): Social modeling for requirements
engineering. MIT Press, Cambridge (2011)
10. Dardenne, A., van Lamsweerde, A., Fickas, S.: Goal directed requirements acquisition. In
Selected Papers of the Sixth Int. Workshop on Software Specification and Design, Elsevier
Science Publishers, (1993), 3--50
11. Bresciani, P., Giorgini, P., Giunchiglia, F., Mylopoulos, J., Perini, A.: TROPOS: An agentoriented software development methodology. International Journal of Autonomous Agents
and Multi-Agent Systems, 8, 3 (2004), 203--236.
12. Giorgini, P., Mylopoulos, J. Nicchiarelli, E., Sebastiani, R.: Formal Reasoning Techniques
for Goal Models. In: Spaccapietra, S. (Ed.) Journal on Data Semantics I, Springer (2004),
1--21.
13. OMG, Business Process Model and Notation, available at:
http://www.omg.org/spec/BPMN/2.0